% Part: sets-functions-relations
% Chapter: sets
% Section: unions-and-intersections

\documentclass[../../../include/open-logic-section]{subfiles}

\begin{document}

\olfileid{sfr}{set}{uni}
\olsection{Verenigingen en doorsneden}

\begin{explain}
In \olref[sfr][set][bas]{sec} hebben we verzamelingen beschreven door te
zeggen aan welke voorwaarde hun elementen moeten voldoen. Dat heet abstractie
en heeft de vorm $\Setabs{x}{\phi(x)}$. In de voorwaarde~$\phi$ mogen
verzamelingen staan die we al hebben beschreven. Neem bijvoorbeeld twee
verzamelingen $A$ en~$B$. In $\Setabs{x}{x \in A \lor x \in B}$ zitten precies
de objecten die !!{element} zijn van $A$ of van~$B$. De !!{element}en van $A$
en~$B$ staan daarin dus bij elkaar. In \olref{fig:union} vormen de
!!{element}en van de twee verzamelingen $A$ en~$B$ samen het gemarkeerde gebied.

\begin{figure}
  \olasset{assets/diagrams/union.tikz}
  \caption{De vereniging $A \cup B$ bestaat uit de !!{element}en van $A$
   samen met de !!{element}en van~$B$.}
  \ollabel{fig:union} 
\end{figure}

Dit samenvoegen van verzamelingen komt vaak van pas. Daarom krijgt deze
bewerking een vaste naam en een symbool.
\end{explain}

\begin{defn}[Vereniging]
De \emph{vereniging} van twee verzamelingen $A$ en $B$ schrijven we als
$A \cup B$. Daarin zitten precies de objecten die !!{element} zijn van $A$,
van $B$ of van allebei.
\[
A \cup B = \Setabs{x}{x \in A \lor x \in B}
\]
\end{defn}

\begin{ex}
Hoe vaak een !!{element} wordt genoemd, maakt voor een verzameling niet uit.
Een !!{element} dat in beide verzamelingen zit, staat daarom maar één keer in
hun vereniging: $\{ a, b, c\} \cup \{ a, 0, 1\} = \{a, b, c, 0, 1\}$.

Als de ene verzameling een deelverzameling van de andere is, is hun vereniging
gewoon de grootste: $\{a, b, c \} \cup \{a \} = \{a, b, c\}$.

Voeg je de lege verzameling toe, dan verandert er niets: $\{a,
b, c \} \cup \emptyset = \{a, b, c \}$.
\end{ex}

\begin{prob}
Laat zien dat als $A \subseteq B$, dan $A \cup B = B$.
\end{prob}

\begin{explain}
Er is ook een bewerking die de ``tegenhanger'' is van verenigen. Die houdt
alleen de !!{element}en over die zowel !!{element} van~$A$ als !!{element}
van~$B$ zijn. Deze bewerking heet de \emph{doorsnede}. Je ziet haar in
\olref{fig:intersection}.
\begin{figure}
  \olasset{assets/diagrams/intersection.tikz}
  \caption{De doorsnede $A \cap B$ bestaat uit de !!{element}en die de twee
    verzamelingen gemeen hebben.}
  \ollabel{fig:intersection}
\end{figure}
\end{explain}

\begin{defn}[Doorsnede]
De \emph{doorsnede} van twee verzamelingen $A$ en $B$ schrijven we als
$A \cap B$. Daarin zitten precies de objecten die zowel !!{element} van $A$
als !!{element} van~$B$ zijn.
\[
A \cap B = \Setabs{x}{x \in A \land x \in B}
\]
Twee verzamelingen heten \emph{disjunct} als hun doorsnede leeg is. Dat
betekent dat ze geen !!{element}en gemeen hebben.
\end{defn}

\begin{ex}
Hebben twee verzamelingen geen !!{element}en gemeen, dan is hun doorsnede leeg:
$\{ a, b, c\} \cap \{ 0, 1\} = \emptyset$.

Hebben ze wel !!{element}en gemeen, dan staan precies die in hun doorsnede:
$\{a, b, c \} \cap \{a, b, d \} = \{a, b\}$.

Als de ene verzameling een deelverzameling van de andere is, is hun doorsnede
de kleinste: $\{a, b, c\} \cap \{a, b\} = \{a, b\}$.

De doorsnede met de lege verzameling is altijd leeg: $\{a, b, c \}
\cap \emptyset = \emptyset$.
\end{ex}

\begin{prob}
Geef een volledig bewijs dat als $A \subseteq B$, dan $A \cap B = A$.
\end{prob}

\begin{explain}
We kunnen ook drie of meer verzamelingen verenigen of doorsnijden. Dat kan op
een mooie, algemene manier. Stop eerst alle verzamelingen waarmee je wilt
werken samen in één verzameling. Hun
vereniging bestaat dan precies uit de objecten die in minstens één
!!{element} van die verzameling zitten. Hun doorsnede bestaat precies uit de
objecten die in ieder !!{element} ervan zitten.
\end{explain}

\begin{defn}
Als de !!{element}en van $A$ zelf verzamelingen zijn, dan bestaat $\bigcup A$
precies uit de !!{element}en die in minstens één !!{element} van~$A$ zitten:
\begin{align*}
\bigcup A & = \Setabs{x}{x \text{ zit in !!a{element} van } A},
\text{ dus:}\\
& = \Setabs{x}{\text{er is een } B \in A
  \text{ met } x \in B}
\end{align*}
\end{defn}

\begin{defn}
Als de !!{element}en van $A$ zelf verzamelingen zijn, dan bestaat $\bigcap A$
precies uit de objecten die in elk !!{element} van~$A$ zitten:
\begin{align*}
\bigcap A & = \Setabs{x}{x \text{ zit in elk !!{element} van } A},
\text{ dus:}\\
 & = \Setabs{x}{\text{voor elke } B \in A, x \in B}
\end{align*}
\end{defn}

\begin{ex}
Neem $A = \{ \{ a, b \}, \{ a, d, e \}, \{ a, d \} \}$.
Dan is $\bigcup A = \{ a, b, d, e \}$ en $\bigcap A = \{ a \}$.
\end{ex}
\begin{prob}
	Laat zien dat als $A$ een verzameling is en $A \in B$, dan
	$A \subseteq \bigcup B$.
\end{prob}

We kunnen hetzelfde doen met een rij verzamelingen $A_1$, $A_2$, \dots
\begin{align*}
\bigcup_i A_i & = \Setabs{x}{x \text{ zit in minstens één van de } A_i}\\
\bigcap_i A_i & = \Setabs{x}{x \text{ zit in elk van de } A_i}.
\end{align*}

Soms gebruiken we een verzameling $I$ als \emph{index}. We bekijken dan
verzamelingen $A_i$, één voor elke $i \in I$. In dat geval mogen we ook dit
schrijven:
\begin{align*}
	\bigcup_{i \in I} A_i & = \bigcup \Setabs{A_i }{i \in I}\\
	\bigcap_{i \in I} A_i & = \bigcap\Setabs{A_i}{i \in I}
\end{align*}

Ten slotte kunnen we kijken naar alles wat wel in~$A$ zit, maar niet in~$B$.
In \olref{difference} zie je dat verschil.

\begin{figure}
  \olasset{assets/diagrams/difference.tikz}
  \caption{Het verschil $A \setminus B$ bestaat uit de !!{element}en die wel
    in $A$ zitten, maar niet in~$B$.}
  \ollabel{difference}
\end{figure}

\begin{defn}[Verschil]
Het \emph{verschil}~$A \setminus B$ bestaat precies uit de !!{element}en die
wel in $A$ zitten, maar niet in~$B$. Dus:
\[
A\setminus B = \Setabs{x}{x\in A \text{ en } x \notin B}.
\]
\end{defn}

\begin{prob}
	Laat zien dat als $A \subsetneq B$, dan $B \setminus A \neq \emptyset$.
\end{prob}
\end{document}
